% Introduction — draft. Contribution bullets are provisional until results exist.
% Do not mention the course assignment or its grading anywhere in the paper.
\section{Introduction}
\label{sec:intro}

Small language models are attractive for domain-specific assistants because they are cheap to fine-tune and serve.
Their weak point is specialized knowledge.
For Korean technical question answering, this weakness takes a particular form.
Technical Korean mixes Korean with English terms and acronyms, so a small model must both recognize the intended technical sense of a term and explain it correctly in Korean.
Training data for this setting is scarce.
The knowledge itself is mostly written in English, and the available domain material, such as lecture slides, is terse, redundant, and partly image-only.

We study this \emph{resource-constrained, domain-specific} setting with a concrete task: answering short Korean questions about the content of a university NLP course with \texttt{Qwen2.5-1.5B-Instruct}.
Answers must be one or two Korean sentences of at most 30 words, generated greedily, and are judged for accuracy, relevance and constraint compliance without a reference answer.

When building fine-tuning data for such a model, two design choices are made almost by default, and their effects are rarely isolated.
The first is \emph{where the knowledge comes from}: only the domain material at hand, or that material enriched with knowledge written by a larger teacher model.
The second is \emph{how technical terms are written}: translated into Korean, or preserved in their original English form within Korean answers.
We ask:
\begin{itemize}
\item[\textbf{RQ1}] Does enriching incomplete domain material with verified teacher-written knowledge improve answer accuracy over material-only data?
\item[\textbf{RQ2}] Does Korean training data that preserves English technical terms improve terminology accuracy and answer accuracy compared with conventionally translated data?
\item[\textbf{RQ3}] Do the two choices interact?
\end{itemize}

We answer these questions with a $2\times2$ factorial study rather than a new method.
We cross knowledge source (lecture-only vs.\ enriched) with terminology policy (translated vs.\ term-preserving) and train all models from the same base model under matched data size, concept coverage and training budget.
\pending{One or two sentences on the main findings, positive or negative, once available.}

Our contributions are \todo{finalize after results}:
\begin{itemize}
\item A reproducible data-construction procedure for domain-specific QA.
The domain material defines the concept scope, and teacher models write canonical knowledge for each concept.
A second model verifies each unit, and material-derived and teacher-written knowledge are kept apart (\S\ref{sec:method}).
\item A controlled $2\times2$ comparison of knowledge enrichment and terminology preservation, including their interaction, with repeated judging, human agreement checks and paired bootstrap intervals (\S\ref{sec:setup}).
\pending{results}
\item An evaluation that separates academic accuracy from terminology accuracy, together with an error analysis of terminology misreading and acronym misinterpretation \pending{results}.
\item A reproducibility note: saving a fine-tuned model with a recent library version can silently switch a greedy evaluation to sampling (Appendix~\ref{app:genconfig}).
\end{itemize}
